% ------------------------------------------------------------
% Page 77
% ------------------------------------------------------------

\begin{center}
    {\Large\bfseries LA VALLÉE DE LA BREZE}
\end{center}

\vspace{0.55cm}

La Brèze se jette dans la Jonte, un peu en amont de Meyrueis, après la piscine et près du
vieux pont désaffecté maintenant. Jusqu’aux années1950 il était dominé par le fameux « arbre
du poivre » que tout le monde connaissait, témoin de bien des jeux d’enfants, de serments
échangés à son ombre comme des heures tragiques que connut notre ville :

\vspace{0.65cm}

\begin{center}
    {\large\bfseries L’ARBRE DU POIVRE}
\end{center}

\vspace{0.45cm}

\begin{center}
\begin{minipage}{0.55\textwidth}
L’arbre du poivre est mort, témoin majestueux\\
Des siècles écoulés. Les flots tumultueux\\
De la Brèze irritée, à son pied tutélaire\\
Ne viendront plus briser l’élan de leur colère.

\vspace{0.35cm}

L’arbre du poivre est mort, il avait vu passer\\
Les héros, qu’un vain roi voulait faire chasser\\
Parce qu’ils priaient Dieu, sans obéir à Rome\\
Sous ses branches, parfois, ils chantaient leurs psaumes,

\vspace{0.35cm}

Ou lorsqu’ils descendaient du roc des protestants\\
Contre son tronc noueux, s’adossaient un moment\\
Disant à demi voix leurs projets téméraires\\
De prêcher au Désert, au risque des galères.

\vspace{0.35cm}

L’arbre du poivre est mort. Les générations,\\
Passent comme un éclair, leurs révolutions\\
Le laissent calme et froid. A son cœur insensible,\\
Les siècles passaient des jours purs et paisibles

\vspace{0.35cm}

Naguère, encore debout dans sa sérénité\\
Il était vigoureux comme en son plein été,\\
Il paraissait doué d’éternelle jeunesse\\
Il semblait à l’abri des coups de la vieillesse

\vspace{0.35cm}

Des atteintes de l’âge il était triomphant\\
Tel que l’avaient toujours connu nos yeux d’enfants\\
Sans cesse rajeuni d’une sève vivante\\
Il élancait au ciel ses branches verdoyantes.
\end{minipage}
\end{center}

\vfill

\begin{flushright}
\textbf{…/…}
\end{flushright}
